% TOP_PROBLEM: {"id": 30005863, "title": "Vacuum Failure of the Third Law of Black-Hole Thermodynamics", "category_id": 16, "category": {"id": 16, "name": "physics", "display_name": "Mathematical Physics", "description": "Problems at the intersection of mathematics and physics.", "slug": "physics", "order_index": 16, "created_at": "2026-07-31T15:26:25.671Z"}, "status": "open", "statusQualification": "The source distinguishes apparent and event horizons; the statement does not require that a single local horizon abruptly change its Kerr parameter. Contextual duplicate 30005862 is excluded."}
% FIRST_POSTED: 2026-10-11
% ENTRY_KIND: statement_only
% SUBMISSION: {"kind": "statement_only", "category": "top_problems", "problem_id": "30005863", "proposal_type": "established", "source": "Analysis, Geometry and Topology of Positive Scalar Curvature Metrics (2024), Oberwolfach Reports, DOI 10.4171/owr/2024/9, pp. 483--566. Original report and mathematical context. Incorporated report material remains under the source CC BY-SA 4.0 terms unless otherwise noted. https://doi.org/10.4171/owr/2024/9", "status_evidence": "Source status retained from the supplied canonical registry. Dated literature assessment (2026-08-22): The claim specifically predicts formation of an exactly extremal four-dimensional Kerr horizon from vacuum gravitational collapse. A 2026 numerical result disproves the third law in five-dimensional vacuum gravity, but it does not prove this four-dimensional Kerr conjecture. Primary report inspected for formulation; no comprehensive current-status literature review is claimed.", "status_checked_on": "2026-10-11", "duplicate_check": "Canonical numeric ID and public code checked against full registry; absent from baseline selected website registry. Source scope compared with adjacent source questions and existing related catalogue entries. No new identity allocated.", "attribution": "Original authors and contributors credited in Analysis, Geometry and Topology of Positive Scalar Curvature Metrics (2024), Oberwolfach Reports, DOI 10.4171/owr/2024/9, pp. 483--566; UnsolvedMath curation credited separately. Definitions and English exposition prepared with Codex assistance; no independent human mathematical verification claimed."}
% SOURCE_STATEMENT: "Construct a vacuum spacetime in which gravitational waves collapse to form a subextremal Kerr apparent horizon and the spacetime later forms an exactly extremal Kerr event horizon."
% Preparation and final self-audit: CONTRIBUTING.md prompts C and D; source/audit artifacts retained in the local physics_statements_2026_10_11 batch.
% Imported from: physics_statements.tex; UnsolvedMath ID 30005863
% !TeX program = lualatex
\documentclass[11pt]{article}
\usepackage[margin=25mm]{geometry}
\usepackage{fontspec}
\setmainfont{Latin Modern Roman}
\usepackage{amsmath,amssymb,mathtools}
\usepackage{unicode-math}
\setmathfont{Latin Modern Math}
\usepackage{xurl,hyperref}
\hypersetup{colorlinks=true,urlcolor=blue}
\setcounter{secnumdepth}{0}
\setlength{\parindent}{0pt}
\setlength{\parskip}{5pt}
\setlength{\emergencystretch}{3em}
\newcommand{\sref}[2]{\hyperlink{source:#1:#2}{[S#2]}}
\newcommand{\cataloguescope}[1]{\paragraph{Recorded target.}#1}
\newcommand{\R}{\mathbb R}
\newcommand{\C}{\mathbb C}
\newcommand{\N}{\mathbb N}
\newcommand{\Z}{\mathbb Z}
\newcommand{\Q}{\mathbb Q}
\newcommand{\F}{\mathbb F}
\newcommand{\eps}{\varepsilon}
\begin{document}
% UnsolvedMath ID 30005863; source catalogue code OWR-14298170-020
\setcounter{section}{30005862}
\section{Vacuum Failure of the Third Law of Black-Hole Thermodynamics}\label{30005863}
{\small UnsolvedMath ID 30005863 · Mathematical Physics}
\subsection{Definitions and mathematical statement}

\paragraph{Shared notation.}$\mathbb N=\{1,2,\ldots\}$; the standard number fields have their usual meanings. All additional conventions are stated locally.


In units $G=c=1$, a Kerr black hole has mass $M$ and angular momentum $J$, with rotation parameter $a=J/M$. Subextremal means $|a|<M$ and extremal means $|a|=M$. An apparent horizon is defined on a spatial slice by trapped-surface geometry; an event horizon is defined globally by visibility from future null infinity. Vacuum means $\operatorname{Ric}(g)=0$.

\paragraph{Catalogue wording retained for traceability.}
Construct a vacuum spacetime in which gravitational waves collapse to form a subextremal Kerr apparent horizon and the spacetime later forms an exactly extremal Kerr event horizon.

\paragraph{Scope and source qualification.}The source distinguishes apparent and event horizons; the statement does not require that a single local horizon abruptly change its Kerr parameter. Contextual duplicate 30005862 is excluded.
\subsection{Short English statement}
Construct gravitational-wave collapse that first produces a subextremal apparent horizon and later an exactly extremal Kerr event horizon.
\subsection{Sources}
\begin{itemize}
\item[\hypertarget{source:30005863:1}{S1}] ULAM.AI and the UnsolvedMath Contributors, UnsolvedMath, record 30005863 (OWR-14298170-020). Catalogue curation under CC BY 4.0; definitions and English exposition added locally.
\url{https://huggingface.co/datasets/ulamai/UnsolvedMath}.
\item[\hypertarget{source:30005863:2}{S2}] Analysis, Geometry and Topology of Positive Scalar Curvature Metrics (2024), Oberwolfach Reports, DOI 10.4171/owr/2024/9, pp. 483--566. Original report and mathematical context. Incorporated report material remains under the source CC BY-SA 4.0 terms unless otherwise noted.
\url{https://doi.org/10.4171/owr/2024/9}.
\end{itemize}
\label{end:30005863}
\end{document}
